\usepackage[margin=1in]{geometry}
\usepackage[unicode]{hyperref}
\usepackage{graphicx}
\usepackage{tikz}
% \usetikzlibrary{arrows.meta,automata,positioning}
\usepackage{import}
\usepackage{apacite}
\usepackage{enumitem}
\usepackage{amsmath,amsfonts,amssymb}


\usepackage{polyglossia}
\setdefaultlanguage{vietnamese}
\setotherlanguage{english}
\addto\captionsvietnamese{%
  \renewcommand{\refname}{Tài liệu tham khảo}%
}


\usepackage{subcaption}
\captionsetup{subrefformat=parens}
\usepackage[skip=4pt]{caption}
\captionsetup*[table]{position=above}

\usepackage{tabularray}
\DefTblrTemplate{firsthead}{caption}{%
  \makebox[\tablewidth]{\parbox{\columnwidth}{%
    \UseTblrTemplate{caption}{normal}%
  }}%
}
\SetTblrTemplate{firsthead}{caption}
\AtBeginDocument{\renewcommand{\tblrcontfootname}{Còn tiếp trang sau}}
\AtBeginDocument{\renewcommand{\tblrcontheadname}{(Tiếp theo)}}
\UseTblrLibrary{varwidth}
\UseTblrLibrary{diagbox}


\usepackage[ruled,vlined]{algorithm2e}
\newcommand{\algorithmname}{Thuật toán}
\SetAlgorithmName{\algorithmname}{\algorithmname}{Danh sách thuật toán}
\DontPrintSemicolon


\usepackage{amsthm}
\renewcommand{\proofname}{Chứng minh}

\theoremstyle{plain} % default
\newtheorem{thm}{Định lý}[section]
\newtheorem{lem}[thm]{Bổ đề}
\newtheorem{prop}[thm]{Mệnh đề}
\newtheorem{cor}[thm]{Hệ quả}

\theoremstyle{definition}
\newtheorem{defn}[thm]{Định nghĩa}
\newtheorem{exmp}[thm]{Ví dụ}

\theoremstyle{remark}
\newtheorem*{rem}{Nhận xét}


\newcommand{\resourcespath}{resources}
\newcommand{\figurespath}{\resourcespath/figures}
\newcommand{\imagespath}{\figurespath}
\graphicspath{{\figurespath}}
\newcommand{\tablespath}{\resourcespath/tables}

\DeclareMathOperator{\expect}{\mathbb{E}}
\DeclareMathOperator{\prob}{\mathbb{P}}
\DeclareMathOperator{\KL}{KL}

\usepackage[toc,xindy={language=vietnamese, codepage=utf8}]{glossaries}
\makeglossaries

\newglossaryentry{filtered}{
  name={được lọc},
  description={filtered}
}

\newglossaryentry{filtration}{
  name={bộ lọc},
  description={filtration}
}

\newglossaryentry{cs}{
  name={dãy tin cậy},
  description={confidence sequence}
}

\newglossaryentry{likelihood-ratio}{
  name={tỉ số khả năng},
  description={likelihood ratio}
}

\newglossaryentry{test}{
  name={kiểm định},
  description={test}
}

\newglossaryentry{st}{
  name={kiểm định tuần tự},
  description={sequential test}
}

\newglossaryentry{acs}{
  name={dãy tin cậy tiệm cận},
  description={asymptotic confidence sequence}
}
